\section{Reproducibility, artifact identity and dependency provenance}
Reproduction needs more than an application version string. Product source,
native harness, reference code, binary, fixtures and environment have separate
identities. The ordinary r4 executable, diagnostic executable and debug native
harness are different artifacts. A successful check on one cannot silently
become an observation of another.

\subsection{Version and input binding}
The current product source is \code{5e59e77}. The native-entry harness was
committed at \code{cb3cd11}; the corrected reference/health code at
\code{b5edb9d}. The exact debug harness executable is unchanged across the two
native attempts. Both protocols snapshot executed helper/native source and
record fingerprints before execution. The evaluator refuses an existing run
directory, changed implementation snapshots or changed protocol/runner.
These controls preserve provenance; they do not independently establish the
correctness of the supplied reference.

The pre-expansion Git evidence checkpoint contains 46 JSON receipts, all
parsed and fingerprint-checked against working files and Git blobs. Reading,
parsing or archiving a receipt is not reproducing its experiment. Historical
line-ending canonicalization preserved the original receipt bytes separately
and did not alter their JSON outcomes. Research source uses LF attributes so
future checkouts can preserve code fingerprints. Raw experimental inputs are
not normalized into a more convenient form.

Each run keeps its protocol, source snapshots, complete synthetic sources,
compiler logs/manifests, PDFs and independent comparison outputs. The invalid
native reference is retained alongside the corrected protocol, rather than
overwritten. The supplied editable-report archive is not a complete application
checkout; running the native harness requires the full recorded repository and
its prepared resources. Reproduction commands and archive scope are in the
appendices.

\subsection{Machine, engines and fonts}
A development environment collector fingerprints explicitly chosen tools,
font files and inputs without compiling or launching the application. The
executed partial snapshot records seven tools, four declared Arial font files
and ten inputs. It was collected after the first reference pilot. It therefore
cannot establish the exact bytes or font resolution used during that earlier
execution. A declared font filename is not proof that no fallback occurred.

RAM, WebView2 and display configuration remain unset in that snapshot. It does
not inventory every shared library, TeX package or background task. Gold builds
for the later native pilot retain recorder output, but complete engine/driver/
OS font binding still needs review before a final freeze. Source documents,
font acquisition and permission to redistribute standalone font files must be
handled separately; the collector does not copy or license those files.

For a final run, seal the source/analysis commits or exact diffs, executables,
worker/DLL, lockfiles, support packages, actual engine/bibliography/font inputs,
OS/CPU/RAM/WebView2/display/DPI/power state, commands, ordering seed and resource
budgets before execution. Document cold/warm and background-load conditions.
Do not generalize one Windows machine's pixels or latency to every installation.

\subsection{The bundled independent verifier}
The packaged verifier uses PDFium 156.0.8076.0 from the pinned pypdfium2 5.14.0
Windows x64 wheel. Historical reports naming PDFium 126 or pypdfium2 4.30 concern
other artifacts. The Rust worker does not require Python at application runtime.
The recorded wheel digest, extracted DLL and packaged DLL match; pinned notices
are included. The index's latest-version observation was made on 4 November
2026, not asserted indefinitely from an old note.

A follow-up checked the wheel's signed PyPI publication against an explicit
expected repository policy, using the attestation mechanism documented by
PyPI~\cite{pypiattest}. The isolated development tool was
\code{pypi-attestations 0.0.31} with \code{sigstore 4.5.1} and default production
trust~\cite{pypitool}. Table~\ref{tab:provenance} records the positive and
negative controls. Verification did not install the wheel or upgrade the app's
runtime, and used no signing or browser login.

\begin{table}[htbp]
\small\begin{tabularx}{\linewidth}{@{}YY@{}}
\toprule
Input / policy & Observed verification result\\
\midrule
Pinned original wheel; expected pypdfium2 repository & Verified signed publication\\
Same basename, valid ZIP with changed archive bytes & Rejected: subject digest mismatch\\
Original wheel; deliberately wrong expected repository & Rejected: repository policy mismatch\\
\bottomrule
\end{tabularx}
\caption{Bounded provenance checks, not application edit-acceptance tests.}
\label{tab:provenance}
\end{table}

The tool was obtained from official PyPI over TLS; its root wheel checksum
matched the index metadata, and installed dependency wheel hashes/versions
were retained. The verifier itself was not independently audited. Seven public
trust metadata files were captured after verification; one \code{root.json}
directory entry was unreadable. That partial post-verification snapshot is
not a presealed trust environment.

Signed publication authenticates an artifact under a policy, not a reproducible
PDFium source-to-binary build, an audit of the publishing workflow or a complete
embedded-component security/license review. An empty package-index advisory
field is not proof of no vulnerabilities. The worker's checked paths, held DLL,
immutable bounded inputs, separate process and time/memory budgets likewise
do not constitute an operating-system sandbox. Nested provenance/build,
vulnerability/version and license review remain public-distribution work.

\subsection{Final freeze rule}
\path{research/evaluation/freeze.json} remains \code{not\_frozen}. The final
fixture directory has no executed held-out cases. Freeze only after the core
workflow and regression gates are stable, labels/oracles have been reviewed,
and final manifests are sealed. If a final evaluation discovers a defect,
retain the failure, fix the development version and start a separately labelled
frozen run. Reusing and tuning the same cases cannot be described as an untouched
final evaluation.
